\documentclass{article}
\usepackage{ctex}        % 中文支持
\usepackage{pgfplots}
\pgfplotsset{compat=1.18}
\usepackage{amsmath}
\usepackage{geometry}
\geometry{a4paper, margin=2.5cm}

\begin{document}

\title{常见函数增长曲线对比}
\author{}
\date{\today}
\maketitle

\section{函数图像}

下图展示了 $\ln x$、$x$、$x^2$ 和 $2^x$ 在算法规模 $n$ 下的增长趋势对比。

\begin{figure}[h!]
\centering
\begin{tikzpicture}
\begin{axis}[
    width=13cm,
    height=9cm,
    xlabel={算法规模 $n$},
    ylabel={运行时间 / 操作次数},
    xmin=0, xmax=10,
    ymin=0, ymax=100,
    grid=both,
    grid style={gray!30},
    legend style={at={(0.02,0.98)}, anchor=north west, cells={anchor=west}},
    legend entries={$\ln n$, $n$, $n^2$, $2^n$},
    samples=200,
    smooth,
    no markers,
    xtick={0,1,2,3,4,5,6,7,8,9,10},
    ytick={0,10,20,30,40,50,60,70,80,90,100},
]

% ln x (绿色) - 从0.1开始避免负无穷
\addplot[domain=0.1:10, thick, color=green!60!black] {ln(x)};

% x (蓝色)
\addplot[domain=0:10, thick, color=blue] {x};

% x^2 (红色)
\addplot[domain=0:10, thick, color=red] {x^2};

% 2^x (紫色) - 只画到6.5以免超出y轴范围
\addplot[domain=0:6.5, thick, color=purple] {2^x};

\end{axis}
\end{tikzpicture}
\caption{$\ln n$、$n$、$n^2$、$2^n$ 的增长曲线对比}
\label{fig:growth}
\end{figure}

$D={a_i|a_i \in ElemSet, i=1, 2, \dots, n, n\ge0}$

\end{document}